% The frustrated instances of Section~\ref{sec:frustrated}, drawn by hand.
% Input by textbook.tex; a sketch and not a rendered figure, so it is outside
% the figure cache and the render cap, which govern generated output only.
\begin{figure}[htbp]
\centering
\begin{tikzpicture}[
    scale=1.05,
    % The shared vocabulary of preamble.tex. The spins are the variables of
    % Figure~\ref{fig:potts:sketch}; only the two labelled states are shaded,
    % and the lattice is drawn without its factors, the bonds standing for
    % them, because what this figure is about is the triangles.
    spin/.style={fgvariable, minimum size=0.48cm, font=\tiny},
    up/.style={spin, fill=black!10},
    down/.style={spin, fill=black!45, text=white},
    bond/.style={fglink},
    agree/.style={draw, very thick, dashed}
]
  % A patch of the triangular lattice: each row offset by half a spacing.
  \foreach \j in {0,1,2} {
    \foreach \i in {0,...,3} {
      \coordinate (n\i\j) at ({\i + 0.5*\j}, {0.87*\j});
    }
  }
  \foreach \j in {0,1,2} {
    \foreach \i in {0,1,2} {
      \draw[bond] (n\i\j) -- ({\i + 1 + 0.5*\j}, {0.87*\j});
    }
  }
  \foreach \j in {0,1} {
    \foreach \i in {0,1,2} {
      \draw[bond] (n\i\j) -- ({\i + 0.5*\j + 0.5}, {0.87*\j + 0.87});
      \draw[bond] ({\i + 1 + 0.5*\j}, {0.87*\j}) -- ({\i + 0.5*\j + 0.5}, {0.87*\j + 0.87});
    }
  }
  \foreach \j in {0,1,2} {
    \foreach \i in {0,...,3} {
      \node[spin] at (n\i\j) {};
    }
  }
  % One triangle: two spins disagree and the third bond cannot.
  \draw[agree] (n00) -- (n01);
  \node[up]   at (n00) {$+$};
  \node[down] at (n10) {$-$};
  \node[up]   at (n01) {$+$};
  \node[anchor=west, font=\footnotesize] at (4.2, 1.30)
    {one agreeing bond per triangle};
  \node[anchor=west, font=\footnotesize] at (4.2, 0.75)
    {$3N$ bonds, $2N$ triangles};
  \node[anchor=west, font=\footnotesize] at (4.2, 0.20)
    {$E_{\min} = |J|\,N$ at zero field};
\end{tikzpicture}
\caption{A patch of the periodic triangular antiferromagnet. Every bond wants
its endpoints to disagree, and a triangle is not two-colourable, so at least
one bond per triangle agrees --- the dashed bond of the labelled triangle,
where $+$ and $-$ are the two states. Counting (triangle, agreeing bond)
pairs over the $2N$ triangles and $3N$ bonds gives
Equation~\ref{eq:triangular-ground-state}, a ground-state energy known at
every size and attained by many configurations. The planted spin glass of
Equation~\ref{eq:planted-energy} is the same picture on a sparse random graph
with a fraction of the couplings inverted. A sketch of the structure and not
a rendering of any instance.}
\label{fig:frustrated:sketch}
\end{figure}
